% ---------------------------------------------------------------------
% 6. Conclusiones y trabajo futuro  (criterio: 10 pts)
% ---------------------------------------------------------------------
\section{Conclusiones y trabajo futuro}
\label{sec:conclusiones}

Este trabajo construyó y validó, sobre 62.043 niños del Observatorio Canguro, un sistema de inteligencia artificial interpretable de dos etapas. Respecto a la pregunta planteada, se concluye que: (i) la calidad del dato es un problema de primer orden y un modelo con contexto detecta el 75\,\% de los errores inyectados, con ganancias grandes donde las reglas de rango fallan; (ii) los riesgos calibrados superan a una regla simple en 13 de 14 desenlaces y son útiles cuando el desenlace depende del estado clínico registrado (oxigenodependencia, desarrollo psicomotor), pero no cuando depende de factores que la base no recoge (no volver al seguimiento); (iii) la comparación con la OMS revela que una de cada cinco tallas está por debajo de $-2$ DE durante el primer año y que la caída de más de 1 DE es una señal más temprana; y (iv) la calibración se desplaza con el tiempo, por lo que la recalibración y la evaluación en silencio son parte del sistema y no un paso opcional. \pendiente{completar con los hallazgos del periodo en silencio retrospectivo, los desenlaces de crecimiento a 12 meses y el despliegue}.

Como trabajo futuro se proponen: un periodo en silencio prospectivo con la siguiente carga semestral de la base y la revisión ciega por neonatólogos; la incorporación de la sede y la fecha de parto, si la Fundación lo autoriza, para vigilar cada sede contra su propio pasado; curvas de prematuros (Fenton o INTERGROWTH-21st) para evaluar el crecimiento antes de las 40 semanas; el análisis de trayectorias y firmas longitudinales de crecimiento, pregunta original de la Fundación, sobre los datos ya depurados; la integración del servicio con la aplicación del Observatorio; y la validación en otros Programas Madre Canguro.
